\centerline{\bf \Huge Первая МатДрака }

\problem{\textbf{1}} Чему равно следующее выражение:  $2^{14}-2^{11}-2^{9}+2^{6}-979$


\problem{\textbf{2}} Сколькими способами можно расставить числа 1 и -1 во всех клетках таблицы 4×4 так, чтобы сумма всех чисел в каждой строке и в каждом столбце была равна 0?


\problem{\textbf{3}} Чему равно следующее выражение:  

$\dfrac{20\dfrac{8}{15}\cdot7.5-54.6:\dfrac{2}{5}}{3\dfrac{13}{21}\cdot8.4-34.4:14\dfrac{1}{3}}+43.75:11\dfrac{2}{3}+24.6:1\dfrac{1}{5}$


\problem{\textbf{4}} 
Найдите наименьшее натуральное $n$ такое, что все 73 дроби $\dfrac{19}{n+21}, \dfrac{20}{n+22},$ 

$\dfrac{21}{n+23}, …, \dfrac{91}{n+93}$ несократимы.


\problem{\textbf{5}} В каждой клетке квадрата 3×3 лежит монета орлом вверх. Какое наименьшее количество монет нужно перевернуть, чтобы в результате не оказалось трех монет, расположенных в одной вертикали, одной горизонтали или одной диагонали и лежащих одинаково (то есть все три орлом вверх или все три решкой вверх)?


\problem{\textbf{6}} 
Треугольная числовая таблица построена таким образом: в вершине, составляющей нулевую строчку, стоит число 0, в $n$-й строчке на левом и правом концах стоит по числу $n$, и между каждыми двумя соседними числами $n$–ой строчки в $(n+1)$–ой стоит их сумма. Найдите сумму всех чисел в сотой строчке.


\problem{\textbf{7}} 
Все цифры трехзначного числа $N$ различны. Из этих цифр составляют всевозможные двузначные числа с неравными цифрами. Найдите все $N$, для которых сумма всех составленных двузначных чисел равна $2N$.


\problem{\textbf{8}} 
Маша пробежала 1 км со средней скоростью 4 м/с. С какой средней скоростью пробежал эту дистанцию Вася, если, стартовав на 25 секунд позже Маши, он финишировал на 25 секунд раньше?


\problem{\textbf{9}} Сколько существует трёхзначных чисел, у которых сумма цифр больше произведения цифр?


\problem{\textbf{10}} Сколько на шахматной доске имеется всевозможных прямоугольников, состоящих из четырёх клеток?

\begin{flushright}
        \scriptsize{}
\end{flushright}

\problem{\textbf{11}} По двум пересекающимся дорогам с равными постоянными скоростями движутся автомобили Ауди и БМВ. Оказалось, что как в 17:00, так и в 18:00 БМВ находился в два раза дальше от перекрестка, чем Ауди. В какое время Ауди мог проехать перекресток? Укажите все возможные варианты.


\problem{\textbf{12}} 
Каждый из 12 человек — рыцарь, всегда говорящий правду, или всегда лгущий лжец. Один из них сказал: «Число лжецов среди нас делится на 1», второй: «Число лжецов среди нас делится на 2», …, 12–ый: «Число лжецов среди нас делится на 12». Сколько среди них может быть рыцарей? Укажите все возможные варианты.


\problem{\textbf{13}} 
Бусы — это кольцо, на которое нанизаны бусины. Бусы можно поворачивать и переворачивать, они от этого не меняются. Сколько различных видов бус можно составить из 10 одинаковых красных бусин и двух одинаковых синих бусин?


\problem{\textbf{14}} 
Найдите количество прямоугольников, составленных из клеток шахматной доски, которые содержат поле C4. (Одна клетка — это тоже прямоугольник.)



\problem{\textbf{15}} 
На шахматной доске стоят ладьи так, что каждая из них бьёт $N$ ладей. При каких целых $N$ это возможно? (Ладья бьёт в каждом направлении только ближайшую ладью.) 


